\chapter{Fluid MoE数学蓝图：约束路由、多速率更新与可验证循环}
\label{chp:fluid_moe_math_blueprint}

本附录把 Fluid MoE 视为一种可供检验的稀疏条件计算方案，而不是智能系统必须服从的物理定律。它所处理的问题很具体：在有限显存、通信、时延和训练样本预算下，系统怎样把任务相关的局部结构约束与全局分布表示联合起来，怎样选择少量专家而不丢失身份、来源和安全边界，又怎样让快速内容更新、较慢路由适配与更慢结构修订保持可追踪的时间关系。我们沿用“形、质、相、耗”等旧名作为工程变量的助记词，但逐一给出状态空间、维度、单位、数值假设和验收指标；这些名称不推出能量守恒、量子坍缩或跨载体整体同构。AgentOS中的 $h(t)$、$E/\Theta$、$\Psi$、TCE、CCM、Boundary、Offloading 与外部记忆 $M_e$ 只是该接口的一种实现，其他智能系统可以采用不同机制，只要提供等价的观测、控制、版本和回退能力。

\section{输入检疫与数值边界：从物理封顶改写为类型化预条件}
\label{sec:appendix_quarantine}

先定义计算对象。设批量大小为 $B$，序列长度为 $N$，模型维数为 $d$，结构通道维数为 $d_f$，内容通道维数为 $d_s$，专家数为 $K$。输入记录写成 $X\in\mathbb{R}^{B\times N\times d}$；它是无量纲的浮点表示，数值大小由上游编码器和归一化约定决定，不是焦耳能量。结构表示 $F\in\mathbb{R}^{B\times N\times d_f}$ 保存任务相关的身份、邻接、角色或位置线索，内容表示 $S\in\mathbb{R}^{B\times N\times d_s}$ 保存可由分布参数读出的属性与上下文。二者并非两个独立本体：同一条记录必须携带来源标识 $r_{bn}$、版本 $v_{bn}$ 与有效掩码 $m_{bn}\in\{0,1\}$，绑定接口 $\operatorname{Bind}(F_{bn},S_{bn};r_{bn},v_{bn})$ 才能声明它们属于同一任务对象。缺少这种身份条件时，即使向量相似，系统也不能把检索结果、用户指令和环境观测任意合并。

所谓“检疫”首先是一组模型预条件。给定可学习投影 $P_f:\mathbb{R}^{d}\to\mathbb{R}^{d_f}$ 与 $P_s:\mathbb{R}^{d}\to\mathbb{R}^{d_s}$，原始通道为 $F^{raw}=P_f(X)$、$S^{raw}=P_s(X)$。若选择定范数结构表示，我们定义
\begin{equation}
F_{bn}=\rho_f\frac{F^{raw}_{bn}}{\max(\lVert F^{raw}_{bn}\rVert_2,\varepsilon_f)},
\qquad m_{bn}=1,
\end{equation}
其中 $\rho_f>0$ 是无量纲目标半径，$\varepsilon_f>0$ 是与浮点精度和归一化尺度共同确定的数值下界。该式只是设计选择：它控制后续点积的尺度，并不证明“逻辑只有方向”。当 $\lVert F^{raw}_{bn}\rVert_2<\varepsilon_f$ 时，直接除法会把噪声放大成貌似确定的方向，因此实现必须走退化分支：保留零向量、读取一个带来源的默认结构先验，或将 $m_{bn}$ 置零并请求上游补充证据。三种策略会产生不同召回率和错误合并率，必须通过消融实验选择，不能用一个微小常数掩盖缺失信息。

内容通道采用均方根约束时，先定义 $q_{bn}=\sqrt{d_s^{-1}\sum_{ell=1}^{d_s}(S^{raw}_{bn\ell})^2}$。给定无量纲区间 $[q_{min},q_{max}]$，可写成
\begin{equation}
S_{bn}=S^{raw}_{bn}\frac{\operatorname{clip}(q_{bn},q_{min},q_{max})}{\max(q_{bn},\varepsilon_s)}.
\end{equation}
这里的 $q$ 是活动强度统计量，不是语义价值、信息熵、推断代价或硬件功耗。若部署需要核算真实资源，应另记显存字节数、通信字节数、浮点运算量、墙钟时延和焦耳能耗组成的资源向量 $c_{bn}$；它们可以与 $q$ 相关，却不能互相替代。下界 $q_{min}$ 也不是越大越好：它可能在弱证据位置制造虚假激活。因而默认实现只裁剪上界，把下界提升交给有明确校准目标的控制器，并报告校准误差、梯度范数分布与异常激活率。

相位变量若确有任务用途，应作为周期特征而非神秘时钟。令 $\theta\in(-\pi,\pi]^{B\times N\times h}$ 表示 $h$ 个路由头的周期坐标，$\nu\in\mathbb{R}^{B\times N\times h}$ 表示每个离散步的相位增量；$\nu$ 的单位是弧度/步，步长为 $\Delta k=1$。边界映射可定义为 $\theta=\pi\tanh(\theta^{raw})$、$\nu=\nu_{max}\tanh(\nu^{raw})$。$\nu_{max}$ 是为数值稳定、采样分辨率和任务时延设定的超参数，不是自然常数，更不应命名为“认知光速”。若系统使用连续时间 $t$，则必须另行声明采样间隔 $\Delta t$，并把频率单位改成弧度/秒；在 $\Delta t$ 变化时仍沿用同一更新式会改变动力学含义。

边界还包括数据与权限，而不只包括数值。掩码 $m$ 应同时排除填充位、越权来源、过期版本和校验失败记录；Boundary 模块保留允许的专家集合 $\mathcal{K}_{bn}\subseteq\{1,\ldots,K\}$，路由器只能在该集合中归一化。对外部记忆 $M_e$ 的读入必须附带来源、时间戳、许可和置信区间，Offloading 结果返回后还要重新绑定原任务身份。一个输入通过范数裁剪，并不意味着它真实、相关或安全；数值检疫、语义校准与权限验证是三个不同关口。

因此，本节的验收不是观察“流形是否稳定”，而是比较可复现指标。至少应记录正常样本上的任务损失与校准误差，极端范数、零向量、长序列和缺失掩码下的有限值比例，身份交换测试中的错误绑定率，越权专家被选中的次数，以及不同精度和步长下的输出偏差。反例同样重要：若无归一化的基线在相同预算下更准确且没有数值故障，定范数投影就没有充分理由；若裁剪降低训练损失却提高罕见事件漏检率，也不能宣称任务成功。只有当预条件在明确分布、指标和资源预算下改善稳健性，并且失败时能够回退，Fluid MoE 的第一层接口才算成立。
